% !TEX root = ../bb_AiRa_Kappaleet.tex

% ö = \%c3\%b6
% ä = \%C3\%A4

\xdef\sectiontitle{\Isectiontitle} %aiheuttama
\TOCrefs

%%%%%%%%%%%%%%%
\begin{frame}[label=1_4_a]%[label=1_4_TOC1]
\frametitle{\texorpdfstring{\culine[\sectiontitleulinecolor]{\sectiontitle}}{\sectiontitle} }
\label{\TOCname}

\vspace{-0.5cm}

\begin{itemize}[leftmargin=0.5cm,itemsep=2mm]
    \item[1.1]<1-> \hyperlink{1_1_a}{\IaSubsectiontitle}
    \item[1.2]<1-> \hyperlink{1_2_a}{\IbSubsectiontitle}	
    \item[1.3]<1-> \hyperlink{1_3_a}{\IcSubsectiontitle}	
    \item[1.4]<1-> \hyperlink{1_4_a}{\CDAlert[red]{\IdSubsectiontitle}}
    \item[1.5]<1-> \hyperlink{1_5_a}{\IeSubsectiontitle}
    \item[1.6]<1-> \hyperlink{1_6_a}{\IfSubsectiontitle}
    \item[1.7]<1-> \hyperlink{1_7_a}{\IgSubsectiontitle}
\end{itemize}

\vspace{-1cm}
\begin{itemize}[leftmargin=8.5cm,itemsep=2mm]
    \item[1.8]<1-> \hyperlink{1_8_a}{\IhSubsectiontitle}
	\item[1.9]<1-> \hyperlink{1_9_a}{\IiSubsectiontitle}
	\item[1.10]<1-> \hyperlink{1_10_a}{\IjSubsectiontitle}
	\item[1.11]<1-> \hyperlink{1_11_a}{\IkSubsectiontitle}
\end{itemize}

%\endofslide 
\end{frame}
%%%%%%%%%%%%%%%

%\xdef\IbSubsectiontitle{Entropy and Temperature}
\xdef\sectiontitle{\IdSubsectiontitle}

%%%%%%%%%%%%%%%
\begin{frame}%[label=1_4_a]
\frametitle{\texorpdfstring{\culine[\sectiontitleulinecolor]{\sectiontitle}}{\sectiontitle} }
\framesubtitle{Boltzmann vs. binomikerroin}
\appear{\xdef\loc{south east}%
\begin{tikzpicture}[remember picture,overlay]% anchor=south
    \node (A) [yshift=-0.25*\figYshift,xshift=\figXshift, anchor=\loc] at (current page.\loc) {\includegraphics[page=1,width=8.2cm]{Figures_booklet/SOM_9_Statistical_mechanics_figures/SOM_Figure_6.png}};%scale=\figscale. % 8cm max hei
\end{tikzpicture}}%
\vspace{-5.5mm}
\begin{itemize}
	\item Todennäköisyyden $P_{n_{i}}$ laskeminen binomikertoimien avulla \appear{(ks. edellinen luku)} \appear{on työlästä suurille systeemeille}\appear{, on järkevämpää käyttää \CDAlert[red]{Boltzmannin jakaumaa}:}
\[ 
 P \textrm{((} E_{n} \textrm{))}  \equal \appear{A} \appear{e^{-E_{n} / k_{B} T} }
 \]
\appear{Kuvassa 6 nähdään että jakauma toimii aika hyvin} \appear{jopa kuvaa pienemmille systeemeille}
\end{itemize}

\vspace{2.5mm}
% 6.5 cm = midway, 10 cm = 3/4 
\begin{minipage}{5.5cm}
\begin{itemize}
	\item Jotta $P \textrm{((} E_{n} \textrm{))} $ voi olla todennäköisyysjakauma\appear{, täytyy se normeerata}\appear{, eli vakio $A$ täytyy löytää/skaalata oikein}
%\[
%\sum P \textrm{((} E_{n} \textrm{))} =\sum_{n} A e^{-E_{n} / k_{B} T}=1 \Rightarrow A=\frac{1}{\sum_{n} e^{-E_{n} / k_{B} T}} \Rightarrow P \textrm{((} E_{n} \textrm{))} =\frac{e^{-E_{n} / k_{B} T}}{\sum_{n} e^{-E_{n} / k_{B} T}}
%\[
%And thus, e.g. average energy can be obtained by
%\[
% < E >\, =\sum_{n} E_{n} P \textrm{((} E_{n} \textrm{))} =\sum_{n} E_{n} \frac{e^{-E_{n} / k_{B} T}}{\sum_{n} e^{-E_{n} / k_{B} T}}=\frac{\sum_{n} E_{n} e^{-E_{n} / k_{B} T}}{\sum_{n} e^{-E_{n} / k_{B} T}}
%\[
\end{itemize}
\end{minipage}

\endofslide 
%\endofsection
\end{frame}
%%%%%%%%%%%%%%%

%%%%%%%%%%%%%%%
\begin{frame}%[label=1_4_a]
\frametitle{\texorpdfstring{\culine[\sectiontitleulinecolor]{\sectiontitle}}{\sectiontitle} }
\framesubtitle{Keskimääräinen energia $< E >$}
\begin{itemize}
\item[] \[
\begin{aligned}
\nsum \appear{P \textrm{((} E_{n} \textrm{))}} &  \equal  \appear{\nsum_{n}} \appear{A} \appear{e^{-E_{n} / k_{B} T} } \equal  \appear{1} \qquad \Rightarrow \qquad \appear{A}  \equal \frac{1}{\nsum_{n} e^{-E_{n} / k_{B} T}} \\
 \appear{\Rightarrow} \qquad \quad \appear{P \textrm{((} E_{n} \textrm{))}} & \equal \frac{e^{-E_{n} / k_{B} T}}{\nsum_{n} e^{-E_{n} / k_{B} T}}
\end{aligned}
\]
\item Nyt, esim. \CDAlert[red]{keskimääräinen energia $< E >\,$}$^*$ voidaan laskea\footnote{{\appear{$^*$\;$< E >\,$ on itseasiassa energian \CDAlert[red]{odotusarvo}{https://fi.wikipedia.org/wiki/Odotusarvo}, joka määritellään yhtälöllä $\mathrm{E}(X) \Equiv \sum_{i=1}^n x_i p_i \,$.} \appear{Opettele~ulkoa.}}}
\[
\begin{aligned}
 < E >\, &  \equal \appear{\nsum_{n}} \appear{E_{n}} \appear{P \textrm{((} E_{n} \textrm{))}} \\
 & \equal \appear{\nsum_{n} E_{n}} \frac{e^{-E_{n} / k_{B} T}}{\nsum_{n} e^{-E_{n} / k_{B} T}} \\
 & \equal \fracc{\nsum_{n} E_{n} e^{-E_{n} / k_{B} T}}{\nsum_{n} e^{-E_{n} / k_{B} T}}
\end{aligned}
\]
\end{itemize}

%\xdef\loc{south east}
%\begin{tikzpicture}[remember picture,overlay] % anchor=south
%    \node (A) [yshift=-0.25*\figYshift,xshift=\figXshift, anchor=\loc] at (current page.\loc) {\includegraphics[page=1,width=8.2cm]{Figures_booklet/SOM_9_Statistical_mechanics_figures/SOM_Figure_6.png}}; %scale=\figscale. % 8cm max hei
%\end{tikzpicture}

\endofslide 
%\endofsection
\end{frame}
%%%%%%%%%%%%%%%

%%%%%%%%%%%%%%%
\begin{frame}
\frametitle{\texorpdfstring{\culine[\sectiontitleulinecolor]{\sectiontitle}}{\sectiontitle} }
\framesubtitle{Jakauman lämpötilariippuvuus}

\begin{itemize}
	\item Einsteinin kiinteälle kappaleelle \appear{päti $E_{n}=n \hbar \omega_{0}$}\appear{, sijoitetaan tämä aiempaan yhtälöön, ja hahmotellaan Boltzmannin jakaumaa (kuva~7)}
	\appear{\xdef\loc{south east}%
\begin{tikzpicture}[remember picture,overlay] % anchor=south
    \node (A) [yshift=-0.25*\figYshift,xshift=\figXshift, anchor=\loc] at (current page.\loc) {\includegraphics[page=1,width=8cm]{Figures_booklet/SOM_9_Statistical_mechanics_figures/SOM_Figure_7.png}}; %scale=\figscale. % 8cm max hei
\end{tikzpicture}}%

\item Kuvassa näkyy Boltzmannin (todennäköisyys)jakauma $P \textrm{((} E_{n} \textrm{))}$ kolmelle eri lämpötilalle  \appear{(tilanteesta $k_{\mathrm{B}} T\!<\!\hbar \omega_{0}$ tilanteeseen $k_{\mathrm{B}} T\!>\!\hbar \omega_{0}$):}
% 6.5 cm = midway, 10 cm = 3/4 
\end{itemize}

\begin{minipage}{6cm}
\begin{itemize}
	\item[] \begin{itemize}[itemsep=2mm]
	\item $T=\Frac{1}{5} \hbar \omega_{0} / k_{B}$
\item $T=\hbar \omega_{0} / k_{B}$
\item $T=5 \hbar \omega_{0} / k_{B}$
\end{itemize}

\item Kylmemmillä lämpötiloilla on selvästi epätodennäköistä että värähtelijät omaisivat useamman energiakvantin \appear{($n\hbar \omega_{0}$)}
\end{itemize}
\end{minipage}
\vspace{2.5mm}


\endofslide 
%\endofsection
\end{frame}
%%%%%%%%%%%%%%%

%%%%%%%%%%%%%%%
\begin{frame}
\frametitle{\texorpdfstring{\culine[\sectiontitleulinecolor]{\sectiontitle}}{\sectiontitle} }
\framesubtitle{Miehitysluku}

\appear{}
\begin{minipage}{7cm}
\begin{itemize}
	\item Meitä kiinnostaa usein tietää montako hiukkasta\appear{/kappaletta}\\  \appear{(värähtelijää)} \appear{on tietyllä tilalla}
		\appear{(omaavat tietyn kvanttiluvun)}

\item Hiukkasten kokonaislukumäärä on edelleen $N$
\end{itemize}
\end{minipage}

\vspace{2.5mm}
\begin{itemize}

\item Määritellään \CDAlert[red]{miehitysluku} $\mathbb{N}$: %$\mathcal{N}$
\[ 
 \mathbb{N} \textrm{((} E_{n} \textrm{))} _{B o l t z m a n n}  \equal  \appear{N P \textrm{((} E_{n} \textrm{))}}  \equal  \appear{N\,} \appear{A e^{-E_{n} / k_{B} T}} \appear{\,, \qquad  < E >\,  \equal \frac{\sum_{n} E_{n} \mathbb{N} \textrm{((} E_{n} \textrm{))} }{\sum_{n} \mathbb{N} \textrm{((} E_{n} \textrm{))} } }
 \]

\item Miehitysluvun $\mathbb{N}$ \appear{avulla  saadaan:}
\[ \bg[2.4]{\text{\bigsymbolsfont(}}\!\! 
\appear{{\footnotesize
\begin{array}{c}
\text {Energialla $E$} \\ \text {olevien hiukkasten} \\ \text {lukumäärä}
\end{array}  }}
 \!\!\bg[2.4]{\text{\bigsymbolsfont)}}   \equal  \appear{\bg[2.4]{\text{\bigsymbolsfont(}\! }
\appear{{\footnotesize
\begin{array}{c}
\mathbb{N} \text {, hiukkasten} \\ \text{lukumäärä tietyllä} \\ 
\text {tilalla E}
\end{array}}}
\!\!\bg[2.4]{\text{\bigsymbolsfont)}}  } \appear{\times} \appear{ \bg[2.4]{\text{\bigsymbolsfont(}}\!\!  
\appear{{\footnotesize 
\begin{array}{c}
\text {Energian E} \\ \text{omaavien tilojen} \\ \text {lukumäärä}
\end{array}}}
\!\!\bg[2.4]{\text{\bigsymbolsfont)}} }
\]
\end{itemize}

\xdef\loc{north east}
\begin{tikzpicture}[remember picture,overlay] % anchor=south
    \node (A) [yshift=\figYshift,xshift=\figXshift, anchor=\loc] at (current page.\loc) {\includegraphics[page=1,width=6.5cm]{Figures_booklet/SOM_9_Statistical_mechanics_figures/SOM_Figure_7.png}}; %scale=\figscale. % 8cm max hei
\end{tikzpicture}

\endofslide 
%\endofsection
\end{frame}
%%%%%%%%%%%%%%%

%%%%%%%%%%%%%%%
\begin{frame}
\frametitle{\texorpdfstring{\culine[\sectiontitleulinecolor]{\sectiontitle}}{\sectiontitle} }
\framesubtitle{Tilojen yli summaaminen vs. integrointi}
\appear{}
\begin{itemize}
	\item Nyt voidaan kirjoittaa: 
	\[
	< E >\,  \equal \fracc{\sum_{n} E_{n} \mathbb{N} \textrm{((} E_{n} \textrm{))} }{\sum_{n} \mathbb{N} \textrm{((} E_{n} \textrm{))} }
	\]
\item Kvanttienergiatilat ovat usein niin lähellä toisiaan, että niiden summaaminen voidaan korvata \CDAlert[red]{integroinnilla}

\item Yleisesti ottaen, kahden vierekkäisen energiatilan todennäköisyyksien suhde on
\[ 
 \Frac{P \textrm{((} E_{n} \textrm{))} }{P \textrm{((} E_{n+1} \textrm{))} }  \equal \fracc{e^{-E_{n} / k_{B} T}}{e^{-E_{n+1} / k_{B} T}} \equal \appear{e^{ \textrm{((} E_{n+1}-E_{n} \textrm{))}  / k_{B} T} }
 \]
\item Jos $E_{n+1}-E_{n}$ on paljon pienempi kuin $k_{B} T$\appear{, suhdeluku on lähellä 1:stä}\appear{, ja summat voidaan muuntaa integraaleiksi}

\end{itemize}

\endofslide 
%\endofsection
\end{frame}
%%%%%%%%%%%%%%%


%%%%%%%%%%%%%%%
\begin{frame}
\frametitle{\texorpdfstring{\culine[\sectiontitleulinecolor]{\sectiontitle}}{\sectiontitle} }
\framesubtitle{Tilojen yli integrointi}
\appear{}
\begin{itemize}
	\item Tämä voidaan näennäisesti tehdä myös raja-arvon käsitteen avulla
\[ 
  < E >\,  \equal \frac{\sum_{n} E_{n} \mathbb{N} \textrm{((} E_{n} \textrm{))} }{\sum_{n} \mathbb{N} \textrm{((} E_{n} \textrm{))} } \qquad \appear{\Rightarrow} \qquad \frac{\int_{n} E \mathbb{N}(E) d n}{\int_{n} \mathbb{N}(E) d n} 
 \]
	\item Tosin summien muuttamisessa integraaleiksi kannattaa olla tarkka
	
\item Kvanttiluku $n$ ei välttämättä ole hyödyllisin integrointimuuttuja\appear{, Energia on parempi muuttuja}

\item Eli meidän pitäisi integroida E:n suhteen:
\[ 
  < E >\,  \equal \frac{\int E \mathbb{N}(E) d n}{\int \mathbb{N}(E) d n} \equal \frac{\int E \mathbb{N}(E) \frac{d n}{d E} \appear{d E}}{\int \mathbb{N}(E) \appear{\Frac{d n}{d E} d E}} \equal \frac{\int E \mathbb{N}(E) D(E) d E}{\int \mathbb{N}(E) D(E) d E} 
 \]

\end{itemize}

\endofslide 
%\endofsection
\end{frame}
%%%%%%%%%%%%%%%

%%%%%%%%%%%%%%%
\begin{frame}
\frametitle{\texorpdfstring{\culine[\sectiontitleulinecolor]{\sectiontitle}}{\sectiontitle} }
\framesubtitle{Tilatiheys}

\begin{itemize}
\item[] %Therefore we should have:
\[ 
  < E >\,  = \Frac{\int E \mathbb{N}(E) d n}{\int \mathbb{N}(E) d n}  = \Frac{\int E \mathbb{N}(E) \Frac{d n}{d E} d E}{\int \mathbb{N}(E) \Frac{d n}{d E} d E} = \Frac{\int E \mathbb{N}(E) D(E) d E}{\int \mathbb{N}(E) D(E) d E} 
 \]
\item Missä on nyt otettu käyttöön uusi suure nimeltään \CDAlert[fuchsia]{energiatilatiheys} \CDAlert[fuchsia]{$D(E)$}\appear{, joka kertoo kuinka tiheästi energiatilat ovat pakkautuneet:}
\[ 
 D(E) \equiv \fracc{\text{\small differentiaalinen tilojen lukumäärä energiavälillä dE lähellä energiaa E}}{d E} \equal \fracc{d n}{d E} 
 \]

%\fontsize{13}{15}\selectfont 
\item Käytännössä suurin osa (jos ei kaikki) keskimääräisten suureiden laskuista \appear{lasketaan integroimalla}
\implies[\normal]{ \CDAlert[red]{Energiatilatiheys $D$ (or DOS)} \appear{\CDAlert[red]{on keskeinen}} \\ \appear{\CDAlert[red]{käsite kiinteän olomuodon fysiikassa!}} }
\implies[\normal]{ \Alert[tunired]{Käsitellään uudelleen/lisää mm.~kiinteän olomuodon fysiikan kurssilla..} }
\end{itemize}

\endofslide 
%\endofsection
\end{frame}
%%%%%%%%%%%%%%%


%%%%%%%%%%%%%%%
\begin{frame}
\frametitle{\texorpdfstring{\culine[\sectiontitleulinecolor]{\sectiontitle}}{\sectiontitle} }
\framesubtitle{Esimerkki}

\appear{}
\begin{itemize}
\item \CDAlert[red]{Harmoninen värähtelijä} on havainnollistava esimerkki näyttää \appear{kuinka energiatilatiheys $D(E)$ voidaan laskea:} 
\[
\begin{aligned}
E_{n}  \equal \appear{n \hbar \omega_{0}} \qquad &\Rightarrow& \qquad
\frac{d n}{d E} &\equal \frac{\text{differentiaalinen tilojen lukumäärä}}{d E} \\
&\Rightarrow& \qquad \appear{d n} &  \equal \frac{1}{\hbar \omega_{0}} \appear{d E} \\
&\Rightarrow &    \frac{d n}{d E} & \equal \frac{1}{\hbar \omega_{0}}  \equal \appear{D(E)}
 \end{aligned}
\]

\item Esimerkki on kieltämättä vähän tylsä\appear{/triviaali}\appear{, kun energiatilatiheys on vakio}\appear{, mutta ainakin tulos on järkevä} \appear{(harmonisen värähtelijän energiatilathan ovat yhtä kaukana toisistaan)}
\end{itemize}

%\xdef\loc{south west}
%\begin{tikzpicture}[remember picture,overlay] % anchor=south
%    \node (A) [yshift=2cm,xshift=-\figXshift, anchor=\loc] at (current page.\loc) {\includegraphics[page=1,width=6.8cm]{Figures_booklet/SOM_9_Statistical_mechanics_figures/SOM_Figure_4_cc.png}}; %scale=\figscale. % 8cm max hei
%\end{tikzpicture}

\endofslide 
%\endofsection
\end{frame}
%%%%%%%%%%%%%%%

%%%%%%%%%%%%%%%
\begin{frame}
\frametitle{\texorpdfstring{\culine[\sectiontitleulinecolor]{\sectiontitle}}{\sectiontitle} }
\framesubtitle{Harmonisten värähtelijöiden energiat}
\appear{}
\begin{itemize}
	\item Identtisille (mutta erotettaville) harmonisille värähtelijöille\appear{, saadaan keskimääräiseksi energiaksi nyt integroitua:}
\[ 
  < E >\,  \equal \fracc{\nint E \mathbb{N}(E) D(E) d E}{\nint \mathbb{N}(E) D(E) d E}  \equal \frac{\nint E N A e^{-E_{n} / k_{B} T} \frac{1}{\hbar \omega_{0}} d E}{\nint N A e^{-E_{n} / k_{B} T} \frac{1}{\hbar \omega_{0}} d E}  \equal \appear{\ldots} \equal \appear{k_{B} T} 
 \]
\item Aiemmin on disktreetille systeemille saatu tulokseksi:
\[ 
  < E >\,  \equal \frac{\hbar \omega_{0}}{e^{\hbar \omega_{0} / k_{B} T}-1} 
 \]
\item Jos on voimassa $\hbar \omega_{0} \ll k_{B} T$\appear{, niin $\appear{x=}\frac{\hbar \omega_{0}}{k_{B} T} \appear{\ll 1}$}\appear{, ja käyttämällä} $\appear{e^{x}} \approx \appear{1+x}$\appear{, saadaan $ \appear{< E >\;} \approx \appear{k_{B} T}$}

\item Molemmat laskutavat antavat saman ennustuksen!
\end{itemize}

%\endofslide 
\endofsection
\end{frame}
%%%%%%%%%%%%%%%

